%% DILARANG EDIT BAGIAN INI
\clearpage
\phantomsection
\addcontentsline{toc}{chapter}{ABSTRAK}
\begin{center}
    \textbf{\large{\MakeUppercase{\judulid}}}\\[0.5cm]
    Oleh\\
    \textbf{\penulis}\\
    NIM. \textbf{\nim}\\[2em]
    \textbf{ABSTRAK}\\[0.5cm]
\end{center}
%% DILARANG EDIT BAGIAN INI

Penelitian ini bertujuan merancang dan membangun aplikasi Point of Sale (POS) berbasis mobile yang dilengkapi fitur prediksi penjualan harian menggunakan metode Long Short-Term Memory (LSTM) untuk membantu pelaku UMKM dalam pengambilan keputusan berbasis data. Sistem dikembangkan menggunakan Flutter sebagai frontend dan Python pada sisi server untuk proses pelatihan serta prediksi model. Data yang digunakan berupa transaksi penjualan harian yang diolah sebagai data deret waktu (time series). Kinerja model dievaluasi menggunakan metrik MAE, RMSE, dan MAPE untuk mengukur tingkat akurasi prediksi. Hasil penelitian diharapkan menghasilkan sistem POS yang tidak hanya mencatat transaksi, tetapi juga mampu memberikan informasi prediksi penjualan guna mendukung perencanaan stok dan strategi usaha secara lebih efektif.\\[0.6cm]

%% DILARANG EDIT BAGIAN INI
\noindent \textbf{Kata Kunci}: Point of Sale (POS), Prediksi Penjualan Harian, Long Short-Term Memory (LSTM), Time Series, Machine Learning, UMKM
%% DILARANG EDIT BAGIAN INI
